\chapter{Introduction}
\label{cap:introducao}

\section{Context and motivation}
\label{sec:motivacao}

The Large Hadron Collider (LHC) at CERN delivers the highest energy proton-proton
collisions ever produced in a laboratory and, with them, an enormous sample of
strongly interacting final states. The majority of these collisions are governed
by Quantum Chromodynamics (QCD), the gauge theory of the strong interaction.
The study and comprehension of these collisions are not only a prerequisite for
suppressing backgrounds in searches for new phenomena but also the structure of
the hadronic final state can be examined and compared to theory directly, at scales
where the coupling is weak enough for perturbation theory to be applied and yet the
dynamics still remain rich.

The most prominent objects in these final states are jets. Jets are collimated sprays
of hadrons that emerge from the fragmentation of high energy quarks and gluons.
A jet is not a fundamental object but instead a definition. A clustering algorithm
specifies which detector inputs are grouped together and a recombination scheme
specifies how momentum is assigned to a result. This defines a jet, from which an
interesting and important question then arises. What does the radiation inside the jets look like?
This domain is called jet substructure, which over the past decade has grown from
a collection of tagging tools into a precision probe of QCD dynamics in its own right. 

In order to organise and look at the internal structure of a jet, one can use the Lund
jet plane \cite{Dreyer:2018nbf}. The Lund jet plane is a two dimensional map of a jet's
branching history on which soft, collinear and non perturbative regions of QCD radiation
occupy distinct territories. This map records where the jet's internal emissions occur, but
it sets aside an intriguing piece of information, the orientation of each splitting
relative to the one before. This orientation is not random. Through the linear polarisation
of intermediate gluons,
successive splittings are correlated in their relative azimuthal angle, producing
a characteristic modulation in the distribution of that angle. This is a genuine quantum
interference effect which survives only because the intermediate gluon is described by
an amplitude rather than a probability, and it is absent from any picture that treats successive
emissions as independent classical branchings.

These collinear spin correlations are well understood theoretically and are implemented in
modern parton shower algorithms, yet establishing them cleanly in collision data has
proved rather difficult. The reason is that the effect carries opposite signs in different
splitting channels. The $g \to q\bar{q}$ and $g \to gg$ channels contribute with opposite
sign, so a measurement that sums all channels sees the two contributions partially cancel.
This interference is not absent from such an inclusive measurement. It is diluted
and the apparent size reflects the accidental flavour composition of the sample rather than the
strength of the underlying physics.

This is the obstacle the present work addresses. A meaningful measurement of the
effect requires knowing which splitting produced each emision. However, the flavour
of an individual splitting inside a jet is not directly observable and must be inferred.
Working entirely on public CERN Open Data \cite{CERNOpenData}, more precisely with
CMS Open Data \cite{CMSOpenDataGuide}, this dissertation builds a reproducible pipeline
from the released samples through Lund declustering to the azimuthal spin observable,
and introduces a flavour selection that isolates the $g \to q\bar{q}$ channel. The
resulting modulation is reported not as a single number but as a function of channel
purity achieved by that selection, a
presentation that separates the physics of the interference from the performance
of the selection and allows the magnitude of the effect to be read off in the
limit of a pure sample.

The Lund declustering angle $\dpsi$ is not the only way to reach this effect.
The same $\cos 2\psi$ modulation also imprints itself on the flow of energy
within a jet, and can be accessed through a three-point energy correlator, which
probes the relative azimuthal orientation of energy deposited by triplets of
particles without reference to any clustering history. The two observables are
complementary: one reads the effect from the reconstructed branching sequence,
the other from the energy flow directly, and each carries a different exposure to
the choices and imperfections of the measurement. This work develops both, 
treating the Lund observable as the primary measurement and the energy correlator
as an independent cross-check on the same underlying physics.

The fact that this analysis rests on public data is not incidental.
The CERN Open Data Portal~\cite{CERNOpenData} has made it possible for work
outside the experimental collaborations to produce meaningful results and to
carry out genuinely new measurements. The present work is intended as a
contribution in that spirit. It builds an end to end pipeline from a single
CMS dataset to the final measurements, establishing along the way a workflow
for accessing, handling and measuring with CMS Open Data.

\section{Objectives}
\label{sec:objetivos}

The main aim of this dissertation is to measure collinear spin correlations in QCD
jets using public LHC data and to establish the flavour selection that such a
measurement requires in order to be meaningful. For the purpose of achieving that goal,
the work is organised around a sequence of objectives. The first is to build a reproducible 
pipeline over CERN Open Data samples capable
of extracting the necessary variables to begin a jet substructure analysis. This means
that the data samples must contain jet constituent information, thus the data format
chosen must retain that level of detail, ruling out the more compact samples in favour
of one that preserves per-constituent quantities.

With that pipeline in place, the second objective is to construct the primary
and secondary Lund jet planes and to verify that they reproduce the known
features of QCD radiation, since this agreement is what establishes that the
declustering is sound before any finer observable is built on top of it. 
Following that, the third objective is to implement the azimuthal observable $\dpsi$
between the hardest primary and secondary branchings, and to study its
distribution both with and without grooming, recovering the inclusive result
before attempting to sharpen it.

The fourth objective is the central one, and the point at which this work departs
from a purely inclusive analysis. Because the spin-interference signal carries
opposite signs in the $g \to q\bar{q}$ and $g \to gg$ channels, an inclusive
measurement dilutes it. Isolating a single channel is therefore not a refinement
but a precondition for measuring the effect at its true strength. This objective
is thus to develop a flavour selection capable of enriching the $g \to q\bar{q}$
channel, and to quantify both how well it performs and how pure a sample it can
reach in principle. The fifth objective builds directly on it. It is to measure
$\langle \cos 2\psi_{12} \rangle$ as a function of the channel purity delivered
by that selection and to establish the value characteristic
of a pure $g \to q\bar{q}$ sample. 

Finally, the work aims to assess the systematic
uncertainties affecting this measurement and to identify the dominant limitations
of the approach, so that the extensions it opens up can be set out on solid ground.

\section{Thesis outline}
\label{sec:estrutura}

Chapter \ref{cap:estado-da-arte} develops the theoretical material on which the
measurement rests, such as the elements of QCD relevant to jet formation, the
definition and reconstruction of jets and the algorithms used to build and
decluster them, the Lund jet plane, and the origin of collinear spin
correlations together with the two observables through which this work accesses
them. Chapter \ref{cap:metodos} turns to the data and the analysis, and covers the CERN
Open Data resources and the CMS detector, the data format and
dataset used, and the pipeline that runs from the released samples through Lund
declustering to the azimuthal observable, the flavour selection, and the
treatment of systematic uncertainties. Chapter \ref{cap:resultados} presents the
validation of the pipeline and the measurement itself, and discusses the results
in the light of the expectations set out in
Chapter~\ref{cap:estado-da-arte}. Finally,
Chapter~\ref{cap:conclusoes} returns to the objectives of
Section~\ref{sec:objetivos}, states what has been established and with what
confidence, and outlines the extensions that the present framework makes
possible.

